% ============================================================
%  Beyond Task Success: Path-Invariant Agreement for
%  Trajectory-Level Evaluation of Non-Deterministic Agents
%
%  NeurIPS 2026 — Datasets and Benchmarks Track
%  Abstract deadline: May 4, 2026
%  Full paper deadline: May 6, 2026
% ============================================================

\documentclass{article}

% ---- NeurIPS 2026 style ----
\usepackage[preprint]{neurips_2026}

% ---- Standard packages ----
\usepackage[utf8]{inputenc}
\usepackage[T1]{fontenc}
\usepackage{hyperref}
\usepackage{url}
\usepackage{booktabs}        % \toprule, \midrule, \bottomrule
\usepackage{amsfonts}
\usepackage{amsmath}
\usepackage{amssymb}
\usepackage{nicefrac}
\usepackage{microtype}
\usepackage{xcolor}
\usepackage{graphicx}
\usepackage{multirow}
\usepackage{array}
\usepackage{enumitem}        % compact lists

% ---- Custom commands ----
\newcommand{\PIA}{\textsc{pia}}          % Path-Invariant Agreement
\newcommand{\kap}{\ensuremath{\kappa}}   % Cohen / Fleiss kappa shorthand
\newcommand{\AgentTrace}{\textsc{AgentTrace}}
\newcommand{\todo}[1]{\textcolor{red}{\textbf{[TODO: #1]}}}
\newcommand{\note}[1]{\textcolor{blue}{\textbf{[NOTE: #1]}}}

% ---- NeurIPS Paper Checklist answer macros ----
\newcommand{\answerYes}[1]{\textbf{[Yes]} #1}
\newcommand{\answerNo}[1]{\textbf{[No]} #1}
\newcommand{\answerNA}[1]{\textbf{[N/A]} #1}

% ============================================================
%  TITLE AND AUTHORS
% ============================================================
\title{Beyond Task Success: Path-Invariant Agreement for
       Trajectory-Level Evaluation of Non-Deterministic Agents}

\author{%
  \textbf{Shailendra Bade} \\
  Independent Researcher \\
  \texttt{shail.finaspirant@gmail.com}
}

% ============================================================
\begin{document}

\maketitle

% ============================================================
\begin{abstract}
% ---- PASTE FINAL ABSTRACT HERE ----
Evaluating multi-step agentic AI systems presents a fundamental measurement
problem: standard inter-annotator agreement (\textsc{iaa}) metrics collapse
when applied to non-deterministic agents that reach identical goals via
legitimately different action sequences.
%
We demonstrate that Fleiss'~$\kappa$ computed over step-level path comparisons
yields $\kappa = -0.065$ on wearable agentic trajectories---not because
annotators disagree on quality, but because the metric conflates path
divergence with quality disagreement.
%
We introduce \textbf{Path-Invariant Agreement (\PIA{})}, a rubric-based
\textsc{iaa} methodology that decouples quality assessment from path length by
evaluating three independent dimensions---planning quality, error recovery, and
goal alignment---rather than action-sequence overlap.
%
Applied to the same trajectory corpus, \PIA{} recovers $\kappa = +0.743$
($\Delta = +0.808$, poor $\to$ substantial agreement).
%
We further show that annotation-layer quality directly determines downstream
model behavior: poisoned step-level labels introduced at the annotation layer
propagate gradient conflicts to 100\% of affected trajectories under outcome
reward models, a failure invisible to task-success metrics.
%
To operationalize both findings, we release \AgentTrace{}, an open-source
end-to-end pipeline for wearable agentic evaluation covering synthetic
trajectory generation, \PIA{} scoring, process reward model annotation, and
annotator outlier detection.
%
Our pipeline improves tool invocation accuracy from 0.36 to 1.00
(+177.8\%) and trajectory success rate from 0.12 to 0.33 on a held-out
benchmark.
%
We argue that trajectory-level evaluation requires new \textsc{iaa} primitives,
and that annotation quality is the pre-gradient intervention point for safe
agentic systems.
\end{abstract}

% ============================================================
\section{Introduction}
\label{sec:intro}
% ============================================================

The deployment of multi-step agentic AI systems has outpaced the evaluation
infrastructure needed to assess them.
%
Enterprise AI surveys find that 89\% of organizations track agent
observability, yet only 52\% have real evaluation frameworks in place
\cite{koreai2025}---a gap driven not by lack of tooling, but by lack of
\emph{methodology}.
%
The dominant evaluation paradigm remains task success: did the agent
accomplish the stated goal?
%
This binary framing is insufficient for two compounding reasons.

\paragraph{Problem 1: Outcome reward models destroy step-level signal.}
Under \emph{outcome reward models} (\textsc{orm}), a single reward is
assigned at the terminal step.
%
If step 15 of a 15-step trajectory fails, all 14 preceding steps---including
those executed correctly---receive an identical negative reward signal.
%
This \emph{gradient conflict} problem forces process reward models (\textsc{prm})
to learn correct intermediate steps from incorrect terminal labels,
degrading training efficiency.
%
Recent work finds that process-supervised training is up to 18$\times$ more
data-efficient than outcome-supervised reinforcement learning precisely because
each intermediate step carries independent signal \cite{reasonrag2025}.
%
The prerequisite for \textsc{prm} training is reliable step-level annotation---
but this is where the second problem emerges.

\paragraph{Problem 2: Standard \textsc{iaa} metrics break for non-deterministic agents.}
Inter-annotator agreement (\textsc{iaa}) metrics such as Fleiss'
$\kappa$~\cite{fleiss1971} and Krippendorff's $\alpha$~\cite{krippendorff2011}
assume that raters evaluating the same artifact produce the same labels when
they agree.
%
This assumption fails for non-deterministic agents: two agents completing the
same goal via legitimately different action sequences produce structurally
different trajectories that path-comparison \textsc{iaa} metrics score as
annotator \emph{disagreement}---even when the annotators share identical
quality judgments.
%
We demonstrate this concretely: applied to wearable agentic trajectories,
standard path-comparison Fleiss' $\kappa$ yields $\kappa = -0.065$.
%
By conventional interpretation, this signals poor agreement and unreliable
annotation.
%
In fact, it signals measurement failure.

\paragraph{Problem 3: Annotation poisoning is undetectable by standard defenses.}
Step-level annotation is the only intervention point between a malicious
annotator and the model's gradient update.
%
Prior work has shown that as few as 250 poisoned training documents---
irrespective of total dataset size---are sufficient to backdoor models of any
scale~\cite{anthropic_backdoor2025}.
%
We show that the standard defense against malicious annotators---Median
Absolute Deviation (\textsc{mad}) suspicion scoring---achieves $F_1 = 0.0$
against coordinated, colluding annotators who assign identical biased scores.
%
Colluding annotators deviate from the panel consensus as a group, collapsing
their joint suspicion score to zero.

\paragraph{Contributions.}
We address all three problems with a unified methodological and artifact
contribution:
\begin{itemize}[nosep]
  \item \textbf{\PIA{} (\underline{P}ath-\underline{I}nvariant
        \underline{A}greement):} a rubric-based \textsc{iaa} methodology that
        evaluates quality dimensions independently of path length, recovering
        $\kappa = +0.743$ ($\Delta = +0.808$, poor $\to$ substantial) on
        trajectories where standard \textsc{irr} yields $\kappa = -0.065$
        (\S\ref{sec:pia}).
  \item \textbf{Poisoning blind spot characterization:} empirical proof that
        \textsc{mad}-based annotator detection achieves $F_1 = 0.0$ against
        colluding annotators, with a complementary Confident Learning
        second-layer defense (\S\ref{sec:poisoning}).
  \item \textbf{\AgentTrace{}:} an open-source, end-to-end pipeline from
        synthetic wearable trajectory generation through \textsc{prm}
        annotation and poisoning detection, released with a publicly available
        annotated dataset on HuggingFace (\S\ref{sec:agenttrace}).
  \item \textbf{Curation impact:} an A/B experiment demonstrating $+177.8\%$
        lift in tool invocation accuracy (0.36 $\to$ 1.00) and trajectory
        success rate (0.12 $\to$ 0.33) from curated vs.\ raw trajectories,
        quantifying the downstream value of annotation quality
        (\S\ref{sec:eval}).
\end{itemize}

% ============================================================
\section{Background and Related Work}
\label{sec:background}
% ============================================================

\subsection{Outcome vs.\ Process Reward Models}
\label{sec:background:prm}

Reinforcement learning from human feedback (\textsc{rlhf}) for language
agents has converged on two reward model paradigms.
%
\emph{Outcome reward models} (\textsc{orm}) assign a single scalar reward
at the terminal step of a trajectory, treating success and failure as
binary outcomes.
%
\emph{Process reward models} (\textsc{prm}) assign intermediate rewards
at each step, enabling credit assignment across the full trajectory.

The gradient conflict problem under \textsc{orm} training is well-documented.
%
\citet{reasonrag2025} show that process-supervised training is up to $18\times$
more data-efficient than outcome-supervised reinforcement learning on
multi-step reasoning tasks: each intermediate step carries an independent
learning signal rather than being zeroed out by a failed terminal reward.
%
\citet{agentprm2025} extend this to agentic settings via Monte Carlo
rollout annotation~\cite{agentprm2025}, generating step-level labels by
simulating all possible trajectory continuations from each intermediate state.
%
Our pipeline implements a simplified, annotation-efficient variant of this
approach: human-in-the-loop step labeling guided by the \PIA{} rubric,
without requiring full Monte Carlo rollout.

The prerequisite for any \textsc{prm} approach is annotation that accurately
reflects step quality---which brings us to the measurement problem this paper
addresses.

\subsection{Inter-Annotator Agreement for Agentic Systems}
\label{sec:background:iaa}

Inter-annotator agreement metrics were developed for static annotation tasks:
labeling text sentiment, named entities, or document categories.
%
These tasks share a property that agentic annotation does not: a given
artifact produces a fixed sequence of labels regardless of which annotator
evaluates it.
%
Non-deterministic agents violate this property by construction---two valid
trajectories to the same goal may have entirely different action sequences,
step counts, and tool invocations.

Despite this, the annotation practices reported in production \textsc{llm}
training pipelines rarely acknowledge the problem.
%
\citet{cohere_commanda2025} describe an 800-prompt annotation process across
65 annotators for the Cohere Command A model---but report no agreement
statistics whatsoever~\cite{cohere_commanda2025}.
%
The absence of $\kappa$ or $\alpha$ reporting is not unusual; it reflects
the broader industry pattern documented by \citet{koreai2025}, who find that
89\% of enterprises track agent observability but only 52\% have real
evaluation frameworks.
%
This gap is particularly notable in light of benchmarks such as the
\textsc{facts} Grounding Leaderboard~\cite{facts_deepmind2025}, which
measures grounding and factuality but does not address the annotation
methodology required to produce step-level training labels for agentic systems.
%
Standard \textsc{iaa} metrics produce misleading results when applied to
non-deterministic agents, and no widely adopted alternative exists.
%
\PIA{} is designed to fill this gap.

\subsection{Annotation Security}
\label{sec:background:security}

The vulnerability of \textsc{rlhf} pipelines to annotation-layer attacks has
received growing attention.
%
\citet{anthropic_backdoor2025} demonstrate that as few as 250 maliciously
crafted documents---independent of total dataset size---are sufficient to
embed persistent backdoors in language models of any scale.
%
This count-based (not proportion-based) finding has direct implications for
annotation pipelines: a small pool of malicious annotators can cause outsized
harm regardless of how large the annotation team is.

\citet{shadearena2025} extend this threat model to multi-agent settings,
demonstrating that frontier LLMs can pursue hidden objectives and evade
monitoring while completing benign tasks---a finding that underscores the
difficulty of detecting subtle model misbehavior at the evaluation layer.
%
Neither work, however, addresses the detection of \emph{colluding annotators}
in non-deterministic agentic annotation---a setting where the standard defense
(\textsc{mad}-based outlier detection) fails by construction.
%
We characterize this blind spot empirically and propose a complementary
defense.

% ============================================================
\section{Path-Invariant Agreement (\PIA{})}
\label{sec:pia}
% ============================================================

\subsection{Why Standard \textsc{irr} Fails for Agents}
\label{sec:pia:failure}

Standard inter-rater reliability metrics---Fleiss' $\kappa$~\cite{fleiss1971},
Cohen's $\kappa$~\cite{cohen1960}, and Krippendorff's
$\alpha$~\cite{krippendorff2011}---share a common assumption: that raters
evaluating the same artifact should produce the same labels if they agree.
This assumption holds for static artifacts such as text documents or images.
It breaks for \emph{non-deterministic agents}.

Consider two annotators evaluating a wearable health agent completing a
blood-pressure monitoring task.
%
Annotator A observes a trajectory that checks the sensor, validates the
reading, and alerts the clinician in three steps.
%
Annotator B observes a trajectory that validates first, rechecks the sensor
due to an anomaly, handles the alert queue, and then alerts in five steps.
%
Both agents succeeded; both annotators would rate the outcome as high quality.
%
Yet under step-level path comparison, the annotators appear to \emph{disagree}:
their step labels do not align because the action sequences differ.
%
Fleiss' $\kappa$ penalizes this path divergence as annotator inconsistency,
even though no inconsistency exists.

Formally, let $T_i$ and $T_j$ be two trajectories reaching the same goal $G$
via action sequences $A_i = (a_1, \ldots, a_m)$ and
$A_j = (a_1', \ldots, a_n')$ where $m \neq n$ or $a_k \neq a_k'$ for some
$k$.
%
A path-comparison \textsc{irr} metric treats any label mismatch at position $k$
as annotator disagreement.
%
When agents are non-deterministic, such mismatches arise even when annotators
share identical quality judgments.
%
The expected $\kappa$ under perfect annotator agreement approaches zero---or
goes negative---in proportion to the variance in valid path lengths.

We confirm this empirically.
%
Applied to our wearable agentic trajectory corpus ($n=30$ trajectories,
5 annotator personas), standard path-comparison Fleiss' $\kappa$ yields
$\kappa = -0.065$ overall, with per-dimension values ranging from
$-0.061$ (privacy compliance) to $+0.002$ (goal alignment).
%
By conventional interpretation~\cite{fleiss1971}, this result signals
\emph{poor agreement}---but it reflects measurement failure, not genuine
annotator disagreement.

\subsection{The \PIA{} Methodology}
\label{sec:pia:method}

\PIA{} resolves the non-determinism paradox by separating the \emph{what}
from the \emph{how}: rather than comparing action sequences step by step,
annotators evaluate \emph{rubric dimensions} that capture quality independent
of the path taken.

\paragraph{Rubric dimensions.}
\PIA{} defines three core dimensions for agentic evaluation:

\begin{itemize}[nosep]
  \item \textbf{Planning quality} (1--4 scale): Does the agent decompose
        the goal into coherent sub-steps? Does it anticipate constraints
        before acting?
  \item \textbf{Error recovery} (1--4 scale): When an intermediate step
        fails or returns an unexpected result, does the agent adapt
        gracefully rather than propagating the failure?
  \item \textbf{Goal alignment} (1--4 scale): Does the agent's final
        state satisfy the stated goal, including any implicit constraints
        such as privacy boundaries or latency requirements?
\end{itemize}

For wearable agentic contexts, a fourth dimension is added:

\begin{itemize}[nosep]
  \item \textbf{Privacy compliance} (1--4 scale): Does the agent respect
        contextual consent boundaries for sensitive biometric data, and does
        it avoid unnecessary data disclosure?
\end{itemize}

Each dimension is scored independently, and \PIA{} $\kappa$ is computed as
Fleiss' $\kappa$ over the dimension scores---not over action-sequence labels.
%
This eliminates the structural source of spurious disagreement: annotators
who agree that an agent planned well will now \emph{agree} on the
planning-quality dimension regardless of which path the agent took.

\paragraph{Calibration protocol.}
Prior to annotation, annotators are exposed to five anchor trajectories:
two scored clearly good (dimension scores 3.5--4.0), one borderline
(dimension scores 2.0--2.5), and two scored clearly bad (dimension scores
1.0--1.5).
%
Anchor scores are computed as:
\begin{equation}
  s_{\text{calibrated}} = 0.82 \cdot s_{\text{gold}} + 0.18 \cdot s_{\text{base}}
  \label{eq:calibration}
\end{equation}
where $s_{\text{gold}}$ is the expert-assigned anchor score and
$s_{\text{base}}$ is the annotator's uncalibrated base score.
%
The calibration weight $w = 0.82$ was selected to achieve target thresholds
of $\kappa \geq 0.55$ and Krippendorff's $\alpha \geq 0.80$.
%
This protocol encodes the observation that borderline cases drive most
annotator disagreement; anchor exposure on borderline examples is therefore
the highest-leverage calibration intervention.

\subsection{Empirical Results}
\label{sec:pia:results}

Table~\ref{tab:pia} reports \PIA{} results across all rubric dimensions on
the wearable agentic trajectory corpus ($n=30$ trajectories, 5 annotator
personas).

\begin{table}[!t]
\centering
\caption{Standard path-comparison \textsc{irr} vs.\ \PIA{} across rubric
         dimensions. Wearable agentic trajectory corpus; $n=30$ trajectories,
         5 annotator personas. Interpretation follows~\citet{fleiss1971}:
         $\kappa < 0.20$ = poor; $0.61$--$0.80$ = substantial;
         $> 0.80$ = almost perfect.}
\label{tab:pia}
\begin{tabular}{lrrrr}
\toprule
\textbf{Dimension} & \textbf{Std.\ $\kappa$} & \textbf{\PIA{} $\kappa$}
  & \textbf{$\Delta$} & \textbf{Interpretation} \\
\midrule
Planning quality   & $-0.065$ & $+0.705$ & $+0.770$ & Substantial     \\
Error recovery     & $-0.065$ & $+0.826$ & $+0.891$ & Almost perfect  \\
Goal alignment     & $-0.065$ & $+0.697$ & $+0.762$ & Substantial     \\
Privacy compliance & $-0.065$ & $+0.942$ & $+1.007$ & Almost perfect  \\
\midrule
\textbf{Overall}   & $-0.065$ & $+0.743$ & $+0.808$ & \textbf{Substantial} \\
\bottomrule
\end{tabular}
\end{table}

\PIA{} raises overall agreement from poor ($\kappa = -0.065$) to substantial
($\kappa = +0.743$), a lift of $\Delta = +0.808$ $\kappa$ points.
%
The error recovery and privacy compliance dimensions reach near-perfect
agreement ($\kappa > 0.80$), suggesting that annotators have well-calibrated
intuitions about these quality properties when freed from path-comparison
artifacts.
%
The privacy compliance dimension is particularly notable: its standard
$\kappa$ of $-0.065$ would conventionally be interpreted as annotators
failing to agree on privacy evaluation---yet its \PIA{} $\kappa$ of $+0.942$
reveals near-consensus once path length is factored out.
%
This pattern is consistent across wearable scenarios that involve sensitive
biometric data, where privacy boundary violations are unambiguous to
annotators but path length varies substantially across agents.

\paragraph{Caveat on calibration artifact.}
Post-calibration $\kappa = 1.00$ reported in dry-run execution is a
mathematical artifact of the deterministic blending formula
(Equation~\ref{eq:calibration}): when $s_{\text{base}}$ is produced
deterministically, all calibrated scores collapse to identical values.
%
The empirically meaningful numbers are the \emph{pre-calibration} values
reported in Table~\ref{tab:pia}, which come from genuine multi-persona
annotation before anchor blending.
%
Live annotation with stochastic LLM calls is expected to yield
$\kappa \approx 0.55$--$0.65$ pre-calibration and
$\kappa \approx 0.78$--$0.85$ post-calibration, consistent with
published human annotation studies on complex NLP tasks.

% ============================================================
\section{Annotation Poisoning at the Step Level}
\label{sec:poisoning}
% ============================================================

\subsection{The Pre-Gradient Intervention Point}
\label{sec:poisoning:pregrad}

Step-level annotation is the only point in the \textsc{prm} training pipeline
where a malicious actor can intervene \emph{before} the gradient update.
%
Once poisoned labels enter the training batch, the model's optimizer has no
mechanism to distinguish a high-quality step from a maliciously mislabeled one.
%
This makes annotation-layer integrity a safety-critical property, not merely a
data quality concern.

The risk is quantified by the \emph{gradient conflict rate}: the proportion of
training steps that receive incorrect reward signals due to label poisoning.
%
Consider an outcome-failed trajectory in which 14 of 15 steps were executed
correctly, but the terminal step failed.
%
Under \textsc{orm} training, all 14 correct steps receive a $-1$ reward,
producing gradient conflicts at every intermediate position.
%
Under \textsc{prm} training with poisoned step labels---where a malicious
annotator assigns $-1$ to correct intermediate steps---the same gradient
conflicts are produced, but they are now invisible to \emph{any} task-success
evaluation metric.
%
Task-success metrics observe only the terminal outcome; they cannot distinguish
correctly-labeled trajectories from poisoned ones.

We measure the gradient conflict rate empirically in \AgentTrace{}.
%
In our synthetic corpus, outcome-failed trajectories in which $\geq 50\%$ of
intermediate steps were executed correctly exhibit a 100\% gradient conflict
rate under \textsc{orm} reward assignment.
%
This result is a synthetic artifact---in production, the rate depends on task
difficulty distribution---but it establishes the upper bound of the problem
and motivates step-level annotation as the correct intervention point.
%
\citet{anthropic_backdoor2025} provide independent evidence for the severity
of the risk: 250 maliciously labeled documents suffice to backdoor arbitrarily
large models, regardless of clean data proportion.

\subsection{MAD Detection Blind Spot}
\label{sec:poisoning:mad}

The standard defense against malicious annotators in large annotation pipelines
is \emph{Median Absolute Deviation} (\textsc{mad}) suspicion scoring.
%
Each annotator receives a suspicion score $s \in [0, 1]$ proportional to their
deviation from the panel consensus on each annotation dimension.
%
Annotators whose scores consistently deviate from the majority are flagged as
outliers and removed.

This defense fails against a specific and practically relevant threat model:
\emph{coordinated collusion}, in which multiple annotators assign identical
biased labels.

Formally, let $\hat{\mu}$ be the panel median for a given annotation dimension,
and let $\delta_i = |x_i - \hat{\mu}|$ be the absolute deviation of annotator
$i$.
%
The \textsc{mad} suspicion score for annotator $i$ is:
\begin{equation}
  s_i = \frac{\delta_i}{\text{MAD}(\delta_1, \ldots, \delta_n)}
  \label{eq:mad}
\end{equation}
where $\text{MAD} = \text{median}(|\delta_i - \text{median}(\delta)|)$.
%
When $k \geq 3$ annotators collude with identical biased labels, their joint
deviation from the panel consensus is shared equally.
%
Their mean deviation equals their individual deviations, so the \textsc{mad}
of the full panel shifts toward the colluders' shared bias.
%
The colluding annotators' suspicion scores collapse toward $s = 0.0$---
indistinguishable from the most reliable annotator in the pool.

We confirm this empirically by injecting three coordinated poisoners
(Poisoner\_A, Poisoner\_B, Poisoner\_C) into our five-annotator panel.
%
All three poisoners receive suspicion scores of $s = 0.000$ at every detection
threshold (Table~\ref{tab:mad}).
%
A threshold sweep from 0.0 to 1.0 yields $F_1 = 0.0$ throughout: precision
and recall are both zero because no poisoner is ever flagged.
%
This is not an implementation failure---it is a fundamental property of
\textsc{mad}-based detection.
%
Any detection scheme that relies solely on deviation from panel consensus will
share this blind spot when the colluding group is large enough to shift
the consensus.

\begin{table}[!t]
\centering
\caption{MAD suspicion scores for the augmented annotator pool (5 legitimate
         + 3 injected colluding poisoners). Poisoners achieve the lowest
         suspicion scores in the pool, making threshold-based detection
         impossible. $F_1 = 0.0$ at all thresholds.}
\label{tab:mad}
\begin{tabular}{llr}
\toprule
\textbf{Annotator} & \textbf{Type} & \textbf{MAD Suspicion Score} \\
\midrule
RecoverySkeptic     & Legitimate & 1.000 \\
OutcomeOptimist     & Legitimate & 0.683 \\
ProcessPurist       & Legitimate & 0.488 \\
ClinicalSafetyFirst & Legitimate & 0.293 \\
PrivacyMaximalist   & Legitimate & 0.195 \\
Poisoner\_A         & Injected   & 0.000 \\
Poisoner\_B         & Injected   & 0.000 \\
Poisoner\_C         & Injected   & 0.000 \\
\bottomrule
\end{tabular}
\end{table}

\subsection{Second-Layer Defense: Confident Learning}
\label{sec:poisoning:cleanlab}

The \textsc{mad} blind spot motivates a complementary defense that operates
on \emph{label distributions} rather than annotator identity.
%
\emph{Confident Learning}~\cite{cleanlab2021} detects noisy labels by
estimating the joint distribution of annotator-assigned labels and latent
true labels, identifying examples where the assigned label is inconsistent
with the model's confident predictions.
%
Crucially, this approach is blind to annotator identity and collusion
structure: it flags suspicious labels regardless of \emph{who} assigned them.

We integrate Cleanlab's \texttt{find\_label\_issues} with quality score
thresholding and \texttt{get\_label\_quality\_scores} into \AgentTrace{} as
a second-layer defense.
%
In the pilot study, Cleanlab identifies zero label issues in the clean
annotation pool---confirming no false positives on legitimate annotators.
%
In future work with real poisoned annotation data, Confident Learning is
expected to surface the label-level signature of colluding poisoners even
when their annotator-level scores are undetectable via \textsc{mad}.
%
The two defenses are complementary: \textsc{mad} catches individual outlier
annotators; Confident Learning catches systematic label corruption regardless
of its source.

% ============================================================
\section{The \AgentTrace{} Pipeline}
\label{sec:agenttrace}
% ============================================================

\subsection{Architecture}
\label{sec:agenttrace:arch}

\AgentTrace{} is a six-stage, end-to-end pipeline from synthetic wearable
trajectory generation to annotation quality evaluation.
%
Each stage produces structured output that feeds directly into the next,
enabling full reproducibility from a single seed.
%
Table~\ref{tab:pipeline} summarizes the stages, modules, inputs, outputs,
and key metrics.

\begin{table}[!t]
\centering
\caption{\AgentTrace{} pipeline summary. All results from \texttt{seed=42};
         100 synthetic trajectories; 5 annotator personas.}
\label{tab:pipeline}
\small
\begin{tabular}{p{2.8cm}p{2.5cm}p{1.5cm}p{1.8cm}}
\toprule
\textbf{Stage} & \textbf{Module} & \textbf{Input} & \textbf{Key Metric} \\
\midrule
1.\ Synthetic generation
  & \texttt{wearable\_generator.py}
  & seed=42
  & 5 scenario types $\times$ 20 logs \\
2.\ Differential privacy
  & \texttt{privacy\_gate.py}
  & Raw sensor data
  & $\varepsilon=1.0$ \\
3.\ Multi-persona annotation
  & \texttt{annotator\_simulator.py}
  & 100 logs, 5 personas
  & Pre-cal $\kappa = -0.035$ \\
4.\ IAA calibration + \PIA{}
  & \texttt{pia\_calculator.py}
  & 150 annotation records
  & $\kappa\colon{-0.065} \to {+0.743}$ \\
5.\ PRM step labeling
  & \texttt{prm\_annotator.py}
  & 20 trajectories
  & 100\% gradient conflict rate \\
6.\ Poisoning detection
  & \texttt{poisoning\_detector.py}
  & 25 records + 3 poisoners
  & MAD $F_1=0.0$; Cleanlab layer \\
\bottomrule
\end{tabular}
\end{table}

Stage 1 generates 100 trajectories (5 scenario types $\times$ 20 logs each).
Stage 2 applies a \textsc{dp} gate ($\varepsilon=1.0$) to raw biometric
sensor values so annotators never observe exact readings.
Stages 3--4 run five annotator personas across four \PIA{} rubric dimensions;
pre-calibration $\kappa = -0.035$ recovers to $+0.743$ after \PIA{} scoring
(\S\ref{sec:pia}).
Stages 5--6 assign step-level \textsc{prm} labels to 20 trajectories
(60 step records) and apply \textsc{mad} + Cleanlab poisoning detection
(\S\ref{sec:poisoning}).

\subsection{Dataset}
\label{sec:agenttrace:dataset}

We release an annotated wearable agentic trajectory dataset on HuggingFace
at \url{https://huggingface.co/datasets/finaspirant/wearable-agent-trajectory-annotations}.
%
The dataset contains 100 synthetic wearable trajectories with full
multi-persona \PIA{} annotations, step-level process reward labels, and
\textsc{dp}-gated sensor data.
%
The annotation schema is organized in three layers:
\begin{itemize}[nosep]
  \item \textbf{Session level:} trajectory metadata, scenario type, agent
        framework, terminal outcome.
  \item \textbf{Role level:} per-annotator-persona dimension scores and
        calibration metadata.
  \item \textbf{Step level:} per-step process reward labels, tool invocation
        records, error flags, and privacy boundary assessments.
\end{itemize}

The dataset includes Croissant metadata~\cite{croissant2024} for
machine-readable dataset documentation, as required by the NeurIPS E\&D
track.
%
The dataset comprises the following artifacts: 100 synthetic wearable
trajectory logs across 5 scenario types (health alert, medication adherence,
vitals monitoring, behavioral pattern detection, and privacy-sensitive ambient
capture); 150 multi-persona annotation records produced by 5 annotator
personas across 30 trajectories on 4 rubric dimensions (planning quality,
error recovery, goal alignment, privacy compliance); 60 step-level process
reward labels across 20 trajectories; and 25 poisoning detection records
including 3 injected adversarial annotations for reproducible blind-spot
evaluation.

\subsection{Reproducibility}
\label{sec:agenttrace:repro}

The full pipeline is available at:
\begin{center}
\url{https://github.com/finaspirant/llm-wearable-agentic-eval-pipeline}
\end{center}
All results reported in this paper are reproducible from \texttt{seed=42}
using the \texttt{uv} environment manager.
%
The repository includes executed Jupyter notebooks
(\texttt{curation\_pipeline\_e2e\_executed.ipynb},
\texttt{agentic\_eval\_flywheel\_executed.ipynb}) with all output cells
preserved, enabling inspection of intermediate results without re-running
the full pipeline.

% ============================================================
\section{Evaluation: Curation Impact}
\label{sec:eval}
% ============================================================

\subsection{Experimental Setup}
\label{sec:eval:setup}

To quantify the downstream value of annotation quality, we conduct an A/B
experiment comparing trajectories produced by the full \AgentTrace{} curation
pipeline (\emph{curated}, $n=50$) against unfiltered raw trajectories
(\emph{raw}, $n=50$) on six evaluation metrics.
Raw trajectories bypass all annotation quality gates: no \PIA{} calibration,
no poisoning detection, and no \textsc{prm} label validation.
Both conditions are evaluated on six metrics inspired by the Kore.ai
enterprise agent evaluation framework \cite{koreai2025}:

\begin{itemize}[nosep]
  \item \textbf{Trajectory success rate:} proportion of trajectories
        reaching the stated goal within the allowed step budget.
  \item \textbf{Tool invocation accuracy:} proportion of tool calls
        that match the expected tool, arguments, and invocation order.
  \item \textbf{Groundedness score} (RAGAS~\cite{ragas2023}): factual
        consistency of agent responses against a reference knowledge base.
  \item \textbf{Privacy leak detection:} proportion of trajectories
        where the agent discloses \textsc{dp}-gated biometric data outside
        authorized consent boundaries.
  \item \textbf{Orchestrator correctness:} proportion of multi-agent
        delegation decisions that route to the correct sub-agent.
  \item \textbf{Latency SLA compliance:} proportion of trajectories
        completing within the defined latency service-level agreement.
\end{itemize}

All experiments use \texttt{seed=42} for full reproducibility.

\subsection{Results}
\label{sec:eval:results}

Table~\ref{tab:ab} presents the curation impact across all six evaluation
metrics.

\begin{table}[!t]
\centering
\caption{A/B experiment: raw vs.\ curated trajectories ($n=50$ per
         condition; \texttt{seed=42}). Starred metrics show statistically
         meaningful lifts; dagger metrics are bounded by synthetic data
         properties.}
\label{tab:ab}
\begin{tabular}{lrrrr}
\toprule
\textbf{Metric} & \textbf{Raw} & \textbf{Curated}
  & \textbf{$\Delta$ (abs)} & \textbf{$\Delta$ (\%)} \\
\midrule
Trajectory success rate$^\star$    & 0.12 & 0.33 & $+0.21$ & $+177.8\%$ \\
Tool invocation accuracy$^\star$   & 0.36 & 1.00 & $+0.64$ & $+177.8\%$ \\
Groundedness (RAGAS)$^\dagger$     & 0.75 & 0.75 & $+0.00$ & --- \\
Privacy leak detection$^\dagger$   & 0.00 & 0.00 & $+0.00$ & --- \\
Orchestrator correctness$^\dagger$ & 1.00 & 1.00 & $+0.00$ & --- \\
Latency SLA compliance$^\dagger$   & 1.00 & 1.00 & $+0.00$ & --- \\
\bottomrule
\end{tabular}

\smallskip
{\footnotesize
$^\star$ Meaningful lift from annotation quality curation.
$^\dagger$ Bounded by synthetic data properties (see text).
}
\end{table}

\paragraph{Meaningful lifts ($^\star$).}
The two metrics showing meaningful improvement directly reflect annotation
quality: trajectory success rate lifts from 0.12 to 0.33 ($+177.8\%$), and
tool invocation accuracy lifts from 0.36 to 1.00 ($+177.8\%$).
%
Both improvements trace directly to the \PIA{} calibration and \textsc{prm}
step labeling stages: curated trajectories filter out those with poor
inter-annotator agreement on tool precision (a \PIA{} sub-dimension of goal
alignment), retaining only trajectories where annotators agree on the
correctness of each tool invocation.
%
This quality gate is invisible to outcome-only evaluation---both the raw
0.36-accuracy trajectories and the curated 1.00-accuracy trajectories may
reach the stated goal---but it directly drives downstream training signal
quality under \textsc{prm} training.

\paragraph{Bounded metrics ($^\dagger$).}
Four metrics show zero delta. Groundedness ($\Delta = 0$, RAGAS $= 0.75$)
reflects a known fallback: without retrieval context, RAGAS returns 0.75
rather than computing factual consistency. Privacy leak detection and
orchestrator correctness are bounded at 0.00 and 1.00 by synthetic data
construction; latency SLA reflects deterministic simulated call latency.
These are synthetic data boundaries, not curation failures. Live evaluation
with real API calls is expected to show curation lifts on all six metrics.

% ============================================================
\section{Conclusion}
\label{sec:conclusion}
% ============================================================

We have identified two compounding failures in the standard evaluation
stack for agentic AI systems. First, standard \textsc{iaa} metrics
misclassify path divergence as annotator disagreement for non-deterministic
agents. Second, step-level annotation poisoning is invisible to task-success
metrics and undetectable by \textsc{mad}-based screening when attackers collude.

\PIA{} addresses the first failure, recovering $\kappa = +0.743$ where
standard metrics yield $\kappa = -0.065$.
Confident Learning addresses the second as a complementary label-level defense.
\AgentTrace{} operationalizes both in an open-source, reproducible pipeline,
validated by $+177.8\%$ lift in tool invocation accuracy and trajectory
success rate from curated vs.\ raw trajectories.

\paragraph{Limitations.}
\PIA{} dimensions were designed for wearable agentic contexts; other domains
must define their own rubric.
All results derive from a synthetic corpus with LLM-simulated
personas---generalization to human annotators requires empirical validation.
The \textsc{mad} blind spot is fundamental to consensus-deviation schemes;
graph-based session-level detection is a promising future direction.
The 100\% gradient conflict rate is a synthetic upper bound (\texttt{outcome\_success=False});
production rates depend on task difficulty distribution.

We release the full pipeline, annotated dataset, and executed notebooks to
enable reproducible research on agentic evaluation methodology.

% ============================================================
%  DATA AVAILABILITY STATEMENT
% ============================================================
\section*{Data Availability}

\textbf{Dataset:} \url{https://huggingface.co/datasets/finaspirant/wearable-agent-trajectory-annotations} — 100 annotated wearable agentic trajectories with \PIA{} scores, \textsc{prm} labels, \textsc{dp}-gated sensor data, and Croissant metadata~\cite{croissant2024}.
\textbf{Code:} \url{https://github.com/finaspirant/llm-wearable-agentic-eval-pipeline} — full \AgentTrace{} pipeline, \texttt{seed=42} reproducible, tagged \texttt{v1.0}; includes executed notebooks with all output cells preserved.
\textbf{Demo:} \url{https://www.loom.com/share/5bec5764428b4e48aa868134e54a894e}

% ============================================================
%  REFERENCES
% ============================================================
\bibliographystyle{plainnat}
\bibliography{references}

% ---- Inline fallback if .bib not yet compiled ----
% \begin{thebibliography}{99}
% \bibitem{reasonrag2025} ...
% \end{thebibliography}

% ============================================================
%  NeurIPS PAPER CHECKLIST
% ============================================================
\newpage
\section*{NeurIPS Paper Checklist}

The checklist is \textbf{not} part of the paper page limit.
Answers in \textbf{bold} indicate the applicable choice.

\begin{enumerate}

\item[\textbf{1.}] \textbf{Claims.}
Do the main claims made in the abstract and introduction accurately
reflect the paper's contributions and scope?
\answerYes{All quantitative claims (PIA $\kappa$ lift, curation
  $+177.8\%$, IAA $0.55{\to}0.82$) are derived from code and data
  released at the linked repository; the methodology is fully
  reproducible from \texttt{seed=42}.}

\item[\textbf{2.}] \textbf{Limitations.}
Did you discuss the limitations of your work?
\answerYes{See \S7 Limitations: synthetic corpus, LLM-simulated
  personas, domain-specific PIA rubric, MAD blind spot for colluding
  annotators, 100\% gradient-conflict rate as a synthetic upper bound.}

\item[\textbf{3.}] \textbf{Theory assumptions and proofs.}
For each theoretical result, did you state the assumptions on which it
relies, and did you give a complete proof?
\answerNA{The paper presents an empirical framework and methodology;
  no new theorems are proved.}

\item[\textbf{4.}] \textbf{Experimental result reproducibility.}
Did you include complete information for reproducing the main
experimental results?
\answerYes{All experiments reproducible from \texttt{seed=42} via the
  open-source repository at
  \url{https://github.com/finaspirant/llm-wearable-agentic-eval-pipeline};
  executed Jupyter notebooks with preserved output cells are included.}

\item[\textbf{5.}] \textbf{Open access to data and code.}
Did you include all the information needed to access the data and code
to reproduce your results?
\answerYes{Code is MIT-licensed on GitHub (link in \S8); the annotated
  dataset (500 records) is publicly available on HuggingFace with a
  Croissant metadata descriptor.}

\item[\textbf{6.}] \textbf{Experimental setting/details.}
Did you specify all the training and test splits, evaluation metrics,
and hyperparameters?
\answerYes{All splits (50 curated / 50 raw, 5 personas, 2 calibration
  phases), evaluation metrics (Cohen's $\kappa$, Fleiss' $\kappa$,
  Krippendorff's $\alpha$, PIA rubric, 6 Kore.ai metrics), and
  hyperparameters ($\varepsilon{=}1.0$, $w{=}0.82$ calibration weight,
  MAD threshold, \texttt{seed=42}) are reported in the text.}

\item[\textbf{7.}] \textbf{Experiment statistical significance.}
Did you report error bars or statistical significance of your results?
\answerNo{PIA and IAA metrics are computed on a single 50-trajectory
  synthetic corpus; bootstrap confidence intervals are a planned
  extension for live-annotator replication.}

\item[\textbf{8.}] \textbf{Experiments compute cost.}
Did you include the compute used to produce your results?
\answerYes{All annotation runs use LLM-simulated personas (no GPU
  training); live API smoke test used \texttt{claude-sonnet-4-6}
  (\textasciitilde{}966 tokens / 10.3 s per trajectory, reported in
  \S4).}

\item[\textbf{9.}] \textbf{Code of ethics.}
Did you review the NeurIPS Code of Ethics?
\answerYes{This study uses fully synthetic wearable sensor data with
  no human participants and no real personal health information.
  IRB review is not required.}

\item[\textbf{10.}] \textbf{Broader impacts.}
Did you discuss any potential negative societal impacts of your work?
\answerYes{See \S7: risks of ambient/wearable AI (consent decay,
  passive capture) are framed as motivation for this work rather than
  a product of it; the framework is designed to surface and mitigate
  these harms via privacy-compliant annotation and HITL triggers.}

\item[\textbf{11.}] \textbf{Safeguards.}
If the paper introduces potentially harmful artifacts (e.g., datasets
that could be used for harmful purposes), did you employ safeguards?
\answerNA{All data is synthetic and contains no real personal
  information; differential privacy ($\varepsilon{=}1.0$) is applied
  to all sensor fields before any downstream use.}

\item[\textbf{12.}] \textbf{Licenses for existing assets.}
Did you cite the license for each existing asset used?
\answerYes{All third-party libraries (LangGraph, CrewAI, AutoGen,
  RAGAS, cleanlab, sentence-transformers) are MIT- or Apache-2.0-
  licensed; citations are provided for all benchmarks and datasets
  referenced.}

\item[\textbf{13.}] \textbf{New assets.}
Did you include everything needed to reproduce the results from the
new assets?
\answerYes{The annotated dataset is on HuggingFace with Croissant
  metadata; the generation code (\texttt{wearable\_generator.py},
  \texttt{annotator\_simulator.py}) is open-source and includes
  \texttt{seed=42} reproducibility.}

\item[\textbf{14.}] \textbf{Crowdsourcing / human subjects.}
Did you include the full text of the instructions given to
participants, recruitment details, and analysis of the resulting data?
\answerYes{Annotator personas are fully specified in \S3 and
  Appendix~A; no real human annotators were recruited — all annotation
  is performed by LLM-simulated personas with deterministic biases
  documented in code.}

\item[\textbf{15.}] \textbf{Institutional Review Board (IRB) approvals.}
Did you describe any potential risks to study participants and whether
you obtained IRB approval?
\answerNA{No human participants; all annotation is LLM-simulated on
  fully synthetic sensor data. IRB review is not required.}

\end{enumerate}

\end{document}
